\documentclass[review]{elsarticle}
\usepackage{lineno}
\usepackage{amssymb}
\usepackage{amsmath}
\usepackage{booktabs}
\usepackage{siunitx}
\usepackage{xcolor}
\usepackage{hyperref}

\journal{Engineering Applications of Artificial Intelligence}

\begin{document}

\begin{frontmatter}

\title{Embedded Artificial Intelligence in Battery Management Systems: Pruning and Quantization for Efficient State-of-Charge and State-of-Health Estimation}

\author[1]{Florian Rzepka\corref{cor1}}
\ead{florian.rzepka@tu-berlin.de}

\author[1]{Julia Kowal}
\ead{julia.kowal@tu-berlin.de}

\cortext[cor1]{Corresponding author.}

\affiliation[1]{organization={Technical University of Berlin, Chair of Electrical Energy Storage Technology},
            addressline={Einsteinufer 11}, 
            city={Berlin},
            postcode={10587}, 
            country={Germany}}

\begin{abstract}
Online battery management requires accurate state-of-charge (SOC) and state-of-health (SOH) estimates within strict memory and timing budgets. To address these constraints, this work investigates how stateful long short-term memory (LSTM) networks with multi-layer perceptron (MLP) heads can be optimised for embedded artificial intelligence (AI), using structured pruning and recurrent-weight quantization to reduce memory and computational costs while retaining estimation accuracy on an STM32H753ZI microcontroller. Using the open \emph{Lithium Iron Phosphate (LFP) Cycle Ageing Dataset}, full-precision models provide the accuracy reference against which the compressed variants are evaluated. Long-horizon replay yields mean absolute error (MAE) values of $2.3$--$2.8$ percentage points for SOC and $0.9$--$1.5$ for SOH, showing that compression largely preserves accuracy. Hardware-in-the-loop benchmarks show that the evaluated 30~\% pruning configuration reduces flash by $40.9$~\% and $45.5$~\% and inference time by $42.9$~\% and $44.0$~\% for the SOC and SOH models, respectively. The analytical complexity model links these savings to pruning fraction and network dimensions. Quantization further reduces flash by $50.2$~\% and $58.8$~\%, but raises runtime in the mixed-precision implementation. Together, the study provides a reproducible reference for resource-efficient embedded battery-state estimation.
\end{abstract}

\begin{keyword}
Embedded artificial intelligence \sep Battery state estimation \sep Long short-term memory network \sep Structured pruning \sep Weight quantization \sep Microcontroller deployment
\end{keyword}

\end{frontmatter}

\section*{Acknowledgements}
This research was funded by the German Federal Ministry for Education and Research (BMBF), grant number 03SF0684A (MG FARM). The authors are responsible for the content. The work was further supported by the German Research Foundation and the Open Access Publication Fund of the Technical University of Berlin. We acknowledge the High Performance Computing Service of the Technical University of Berlin for providing the computational resources used in the pruning and quantization experiments. The authors thank Nils Strassenburg (Hasso-Plattner-Institut, Potsdam) for helpful discussions on pruning and quantization, Mano Schmitz for the collaboration on the dataset creation, and Max M\"onikes and Pavel Aleinikau (Technical University of Berlin) for their collaboration on the BMS hardware.

\section*{CRediT authorship contribution statement}
\textbf{Florian Rzepka}: Conceptualization, methodology, software, validation, formal analysis, investigation, data curation, writing--original draft preparation, writing--review and editing, visualization.
\textbf{Julia Kowal}: Resources, supervision, project administration, funding acquisition, writing--review and editing.

\section*{Declaration of competing interest}
The authors declare that they have no known competing financial interests or personal relationships that could have appeared to influence the work reported in this paper.

\section*{Declaration of generative AI and AI-assisted technologies in the manuscript preparation process}
During the preparation of this work the authors used ChatGPT, DeepL, and Grammarly to improve the linguistic quality and style of the manuscript. After using these tools, the authors reviewed and edited the content as needed and take full responsibility for the content of the published article.

\end{document}
